\scanpage{42}\mspage{35}
% archive sheet number: 41
\begingroup\makeatletter\tagsleft@true\makeatother
Les structures de Quillen faibles \deleted{pour} les $(M(I),W(I))$ seront déduites d'une structure \deleted{\unsure{de} $\unsure{Qf}$ pour} $(M,W)$, de la façon suivante : elles seront toutes ``\uline{closes à droite}'' i.e.
\[
 \deleted{\unsure{M}}\ F(I)=(TC(I))_{*},\quad TF(I)=(C(I))_{*},\tag*{(\unsure{1})}
\]
\endgroup
en particulier (pour $I=e$)
\[
 F=(TC)_{*},\quad TF=C_{*}
\]
ce qui implique déjà (puisque $C\subset\operatorname{Cof}_{W},TC\subset W^{\unsure{\mathrm{univ}|\mathrm{uni}}}$)
\[
 (\operatorname{Cof}_{W})_{*}\subset W_{\deleted{\unsure{\mathrm{univ}}}},\quad(W^{\mathrm{univ}})_{*}\subset\operatorname{Fib}_{W},\tag*{(\unsure{2})}
\]
\inserted{(d'où $(\operatorname{Cof}_{W})_{*}\subset W_{\mathrm{univ}}$).}
En vertu de ça, pour définir les structures $Qf$ \deleted{pour les} $(M(I),W(I))$, il suffit de définir $C(I),TC(I)$, et $=$, comme on veut l'exactitude à gauche i.e.
\[
 TC(I)=C(I)\cap W(I)\tag*{(\unsure{3})}
\]
(en particulier
\[
 TC=C\cap W\qquad),\tag*{(\unsure{4})}
\]
il reste à définir $C$. On \deleted{\unsure{prendra}} \inserted{définira} $C(I)$ ``argument \deleted{\unsure{et}} par argument'', i.e. pour $f\in\operatorname{Fl}(M(I))$
\[
 f\in C(I)\Longleftrightarrow\alpha_{s}^{*}(f)\in C\quad\text{pour tt }s\in\operatorname{Ob}I,\tag*{(5)}
\]
où $\alpha_{s}:\deleted{\unsure{\{s\}}}e\longrightarrow I$ défini par $\alpha_{s}(e)=s$.
Il en résulte d'ailleurs aussitôt, \deleted{par} (3) \deleted{\unsure{et}}
\[
 f\in TC(I)\Longleftrightarrow\alpha_{s}^{*}(f)\in TC\quad\text{pour tt }s\in\operatorname{Ob}I.\tag*{(6)}
\]
